\documentclass[10pt,a4paper]{amsart}
\usepackage{amsmath,amssymb,amsthm,graphicx}
\usepackage{verbatim}
\newenvironment{reference}{\comment}{\endcomment}
%\newenvironment{reference}{}{}
\newenvironment{slogan}{\comment}{\endcomment}
\newenvironment{history}{\comment}{\endcomment}

% For commutative diagrams we use Xy-pic
\usepackage[all]{xy}

% We use 2cell for 2-commutative diagrams.
\xyoption{2cell}
% Theorem environments.
%
\theoremstyle{plain}
\newtheorem{theorem}[subsection]{Theorem}
\newtheorem{proposition}[subsection]{Proposition}
\newtheorem{lemma}[subsection]{Lemma}

\theoremstyle{definition}
\newtheorem{definition}[subsection]{Definition}
\newtheorem{example}[subsection]{Example}
\newtheorem{exercise}[subsection]{Exercise}
\newtheorem{situation}[subsection]{Situation}

\theoremstyle{remark}
\newtheorem{remark}[subsection]{Remark}
\newtheorem{remarks}[subsection]{Remarks}

\numberwithin{equation}{subsection}

% Macros
%
\def\lim{\mathop{\mathrm{lim}}\nolimits}
\def\colim{\mathop{\mathrm{colim}}\nolimits}
\def\Spec{\mathop{\mathrm{Spec}}}
\def\Hom{\mathop{\mathrm{Hom}}\nolimits}
\def\Ext{\mathop{\mathrm{Ext}}\nolimits}
\def\SheafHom{\mathop{\mathcal{H}\!\mathit{om}}\nolimits}
\def\SheafExt{\mathop{\mathcal{E}\!\mathit{xt}}\nolimits}
\def\Sch{\mathit{Sch}}
\def\Mor{\mathop{\mathrm{Mor}}\nolimits}
\def\Ob{\mathop{\mathrm{Ob}}\nolimits}
\def\Sh{\mathop{\mathit{Sh}}\nolimits}
\def\NL{\mathop{N\!L}\nolimits}
\def\CH{\mathop{\mathrm{CH}}\nolimits}
\def\proetale{{pro\text{-}\acute{e}tale}}
\def\etale{{\acute{e}tale}}
\def\QCoh{\mathit{QCoh}}
\def\Ker{\mathop{\mathrm{Ker}}}
\def\Im{\mathop{\mathrm{Im}}}
\def\Coker{\mathop{\mathrm{Coker}}}
\def\Coim{\mathop{\mathrm{Coim}}}

% Boxtimes
%
\DeclareMathSymbol{\boxtimes}{\mathbin}{AMSa}{"02}

%
% Macros for moduli stacks/spaces
%
\def\QCohstack{\mathcal{QC}\!\mathit{oh}}
\def\Cohstack{\mathcal{C}\!\mathit{oh}}
\def\Spacesstack{\mathcal{S}\!\mathit{paces}}
\def\Quotfunctor{\mathrm{Quot}}
\def\Hilbfunctor{\mathrm{Hilb}}
\def\Curvesstack{\mathcal{C}\!\mathit{urves}}
\def\Polarizedstack{\mathcal{P}\!\mathit{olarized}}
\def\Complexesstack{\mathcal{C}\!\mathit{omplexes}}
% \Pic is the operator that assigns to X its picard group, usage \Pic(X)
% \Picardstack_{X/B} denotes the Picard stack of X over B
% \Picardfunctor_{X/B} denotes the Picard functor of X over B
\def\Pic{\mathop{\mathrm{Pic}}\nolimits}
\def\Picardstack{\mathcal{P}\!\mathit{ic}}
\def\Picardfunctor{\mathrm{Pic}}
\def\Deformationcategory{\mathcal{D}\!\mathit{ef}}
\usepackage{cleveref}
\crefname{originallemma}{original lemma}{original lemmas}
\Crefname{originallemma}{Original lemma}{Original lemmas}
\makeatletter
\newlabel{algebra-lemma-pullback-module}{{02M1}{source}{}{sourcepointer-1}{}}
\newlabel{algebra-lemma-pullback-module@cref}{{[originallemma][0][]02M1}{[1][0][]source}{}{}{}}
\makeatother
\crefname{originalsection}{original section}{original sections}
\Crefname{originalsection}{Original section}{Original sections}
\makeatletter
\newlabel{decent-spaces-section-residue-fields-henselian-local-rings}{{0EMV}{source}{}{sourcepointer-2}{}}
\newlabel{decent-spaces-section-residue-fields-henselian-local-rings@cref}{{[originalsection][0][]0EMV}{[1][0][]source}{}{}{}}
\makeatother
\crefname{originalreference}{original reference}{original references}
\Crefname{originalreference}{Original reference}{Original references}
\makeatletter
\newlabel{situation-setup}{{02QL}{source}{}{sourcepointer-3}{}}
\newlabel{situation-setup@cref}{{[originalreference][0][]02QL}{[1][0][]source}{}{}{}}
\makeatother
\crefname{originaldefinition}{original definition}{original definitions}
\Crefname{originaldefinition}{Original definition}{Original definitions}
\makeatletter
\newlabel{spaces-divisors-definition-zero-scheme-s}{{0845}{source}{}{sourcepointer-4}{}}
\newlabel{spaces-divisors-definition-zero-scheme-s@cref}{{[originaldefinition][0][]0845}{[1][0][]source}{}{}{}}
\makeatother
\crefname{originaldefinition}{original definition}{original definitions}
\Crefname{originaldefinition}{Original definition}{Original definitions}
\makeatletter
\newlabel{spaces-over-fields-definition-order-vanishing-meromorphic}{{0EPR}{source}{}{sourcepointer-5}{}}
\newlabel{spaces-over-fields-definition-order-vanishing-meromorphic@cref}{{[originaldefinition][0][]0EPR}{[1][0][]source}{}{}{}}
\makeatother
\crefname{originalsection}{original section}{original sections}
\Crefname{originalsection}{Original section}{Original sections}
\makeatletter
\newlabel{spaces-over-fields-section-c1}{{0ENV}{source}{}{sourcepointer-6}{}}
\newlabel{spaces-over-fields-section-c1@cref}{{[originalsection][0][]0ENV}{[1][0][]source}{}{}{}}
\makeatother
\title{The divisor of a nonzero global section equals the class of its zero scheme and computes the first Chern class}
\author{The Stacks Project Authors}
\date{stacks-project:spaces-chow.tex:local-snapshot-20261012; GNU FDL 1.2 or later}
\begin{document}
\maketitle
\noindent\textit{Editorial note.} All delivered mathematical text is the author's own, copied from the pinned Stacks Project chapter file spaces-chow.tex; no statement or proof step was rewritten, reordered or removed.\par\medskip
\noindent\textit{Editorial note.} The counted proof is the author's complete argument: an elementary etale neighbourhood trivialising the sheaf, the identification of the zero scheme with the zero scheme of a local equation, and a length comparison through the henselization\par\medskip
\noindent\textit{Editorial note.} The label lemma-compute-c1 names the counted statement, and the proof environment immediately following it in the source is delivered verbatim as its proof.\par\medskip
\noindent\textit{Editorial note.} The counted statement is carried by the author's own lemma environment in the chapter source.\par\medskip
\noindent\textit{Allowed standard prerequisites (editorial).}\begin{itemize}
\item The zero scheme of a nonzero global section of an invertible sheaf and the divisor of a section, together with the order of vanishing of a section along a prime divisor
\item Elementary etale neighbourhoods and the identification of the henselizations of the local rings of a point and of the corresponding point of the neighbourhood
\item Flat base change of lengths along the map from a local ring to its henselization, and the identification of the coefficients of a cycle by comparing them at generic points
\end{itemize}\medskip
\ifdefined\hypertarget\hypertarget{sourcepointer-1}{}\fi
\noindent\textit{Source reference.} Original lemma 02M1: sources/stacks-chow/tags.txt, source lines 3428--3428. Prerequisite pointer invoked at proof line(s) 2375; it names a standard fact of a sibling chapter that the counted argument uses, and every step of the counted proof built on it is retained verbatim in the delivered body.\par\medskip
\ifdefined\hypertarget\hypertarget{sourcepointer-2}{}\fi
\noindent\textit{Source reference.} Original section 0EMV: sources/stacks-chow/tags.txt, source lines 18320--18320. Prerequisite pointer invoked at proof line(s) 2350; it names a standard fact of a sibling chapter that the counted argument uses, and every step of the counted proof built on it is retained verbatim in the delivered body.\par\medskip
\ifdefined\hypertarget\hypertarget{sourcepointer-3}{}\fi
\noindent\textit{Source reference.} Original reference 02QL: sources/stacks-spaces-chow/pinned-original-source.tex, source lines 264--264. Prerequisite pointer invoked at proof line(s) 2329; it names a standard fact of the same chapter that the counted argument uses, and every step of the counted proof built on it is retained verbatim in the delivered body.\par\medskip
\ifdefined\hypertarget\hypertarget{sourcepointer-4}{}\fi
\noindent\textit{Source reference.} Original definition 0845: sources/stacks-chow/tags.txt, source lines 10310--10310. Prerequisite pointer invoked at proof line(s) 2357; it names a standard fact of a sibling chapter that the counted argument uses, and every step of the counted proof built on it is retained verbatim in the delivered body.\par\medskip
\ifdefined\hypertarget\hypertarget{sourcepointer-5}{}\fi
\noindent\textit{Source reference.} Original definition 0EPR: sources/stacks-chow/tags.txt, source lines 18386--18386. Prerequisite pointer invoked at proof line(s) 2378; it names a standard fact of a sibling chapter that the counted argument uses, and every step of the counted proof built on it is retained verbatim in the delivered body.\par\medskip
\ifdefined\hypertarget\hypertarget{sourcepointer-6}{}\fi
\noindent\textit{Source reference.} Original section 0ENV: sources/stacks-chow/tags.txt, source lines 18355--18355. Prerequisite pointer invoked at proof line(s) 2358; it names a standard fact of a sibling chapter that the counted argument uses, and every step of the counted proof built on it is retained verbatim in the delivered body.\par\medskip
\begin{lemma}
\label{lemma-compute-c1}
In Situation \ref{situation-setup} let $X/B$ be good.
Assume $X$ is integral and $n = \dim_\delta(X)$.
Let $\mathcal{L}$ be an invertible $\mathcal{O}_X$-module.
Let $s \in \Gamma(X, \mathcal{L})$ be a nonzero global section.
Then
$$
\text{div}_\mathcal{L}(s) = [Z(s)]_{n - 1}
$$
in $Z_{n - 1}(X)$ and
$$
c_1(\mathcal{L}) \cap [X] = [Z(s)]_{n - 1}
$$
in $\CH_{n - 1}(X)$.
\end{lemma}

\begin{proof}
Let $Z \subset X$ be an integral closed subspace of $\delta$-dimension $n - 1$.
Let $\xi \in |Z|$ be its generic point. To prove the first equality
we compare the coefficients of $Z$ on both sides. Choose an elementary
\'etale neighbourhood $(U, u) \to (X, \xi)$, see
Decent Spaces, Section
\ref{decent-spaces-section-residue-fields-henselian-local-rings}
and recall that $\mathcal{O}_{X, \xi}^h = \mathcal{O}_{U, u}^h$
in this case. After replacing $U$ by an open neighbourhood
of $u$ we may assume there is a trivializing section
$s_U$ of $\mathcal{L}|_U$. Write $s|_U = f s_U$ for some
$f \in \Gamma(U, \mathcal{O}_U)$. Then $Z \times_X U$ is
equal to $V(f)$ as a closed subscheme of $U$, see
Divisors on Spaces, Definition \ref{spaces-divisors-definition-zero-scheme-s}.
As in Spaces over Fields, Section \ref{spaces-over-fields-section-c1}
denote $\mathcal{L}_\xi$ the pullback of $\mathcal{L}$
under the canonical morphism
$c_\xi : \Spec(\mathcal{O}_{X, \xi}^h) \to X$.
Denote $s_\xi$ the pullback of $s_U$; it is a trivialization of
$\mathcal{L}_\xi$. Then we see that
$c_\xi^*(s) = fs_\xi$. The coefficient of $Z$ in $[Z(s)]_{n - 1}$
is by definition
$$
\text{length}_{\mathcal{O}_{U, u}}(\mathcal{O}_{U, u}/f\mathcal{O}_{U, u})
$$
Since $\mathcal{O}_{U, u} \to \mathcal{O}_{X, \xi}^h$
is flat and identifies residue fields this is equal to
$$
\text{length}_{\mathcal{O}_{X, \xi}^h}
(\mathcal{O}_{X, \xi}^h/f\mathcal{O}_{X, \xi}^h)
$$
by Algebra, Lemma \ref{algebra-lemma-pullback-module}.
This final quantity is equal to $\text{ord}_{Z, \mathcal{L}}(s)$ by
Spaces over Fields, Definition
\ref{spaces-over-fields-definition-order-vanishing-meromorphic}, i.e.,
to the coefficient of $Z$ in
$\text{div}_\mathcal{L}(s)$
as desired.
\end{proof}

\end{document}
